\documentclass[aspectratio=169,11pt]{beamer}

\usetheme{Madrid}
\usecolortheme{default}
\usefonttheme{professionalfonts}
\setbeamertemplate{navigation symbols}{}
\setbeamertemplate{footline}[frame number]

\usepackage{amsmath, amssymb, booktabs}

\newcommand{\E}{\mathbb{E}}
\newcommand{\1}{\mathbf{1}}
\newcommand{\plim}{\operatorname{plim}}

\title{Curating Maps And Spatial Data}
\subtitle{Empirical Methods --- Week 15}
\author{Teaching deck draft}
\date{}

\begin{document}

\begin{frame}
  \titlepage
\end{frame}

\begin{frame}{Week 15 question}
\begin{itemize}
    \item How do researchers turn places, borders, distances, and exposure measures into empirical objects?
    \item Spatial curation comes before spatial causal inference.
    \item Geography choice, exposure definitions, crosswalks, and travel costs shape the estimand.
    \item A map is not evidence until the measurement object is justified and auditable.
\end{itemize}
\end{frame}

\begin{frame}{Spatial data are economic objects}
\small
\[
c_{ijm}=\tau_m t_{ijm}+\pi_m d_{ijm}+\kappa_m
\]

\[
E_i=\sum_j w_{ij}s_j
\]

\begin{itemize}
    \item \(c_{ijm}\): generalized access or commuting cost by mode \(m\).
    \item \(s_j\): shock, opportunity, amenity, policy, or disamenity at place \(j\).
    \item \(w_{ij}\): the researcher's exposure technology.
    \item The unit and weights define the object before any estimator is run.
\end{itemize}
\end{frame}

\begin{frame}{Spatial objects}
\small
\begin{tabular}{p{0.18\linewidth} p{0.32\linewidth} p{0.38\linewidth}}
\toprule
Object & Best for & Main risk \\
\midrule
Point & addresses, plants, schools, stops & geocoding error, confidentiality, false precision \\
Polygon & tracts, counties, commuting zones & MAUP, ecological fallacy, arbitrary borders \\
Raster & pollution, heat, night lights, terrain & resolution, temporal match, aggregation \\
Network & roads, transit, commute links & mode assumptions, changing schedules, reproducibility \\
Panel geography & units over time & boundary changes, crosswalk error, support \\
\bottomrule
\end{tabular}
\end{frame}

\begin{frame}{Geographies and estimands}
\small
\begin{tabular}{p{0.25\linewidth} p{0.31\linewidth} p{0.32\linewidth}}
\toprule
Research object & Plausible geography & Estimand discipline \\
\midrule
Local labor-market shock & commuting zone, county & Is the market integrated by commuting? \\
Neighborhood exposure & tract, block, residence history & Does timing match the relevant exposure window? \\
Job access & travel-time isochrone, reachable jobs & Are jobs reachable by the relevant workers? \\
Policy jurisdiction & city, school zone, county & Does treatment follow the legal boundary? \\
Market access & network, trade-cost matrix & Are routes and costs measured for the period? \\
\bottomrule
\end{tabular}

\medskip
The same raw locations can answer different questions under different geographies.
\end{frame}

\begin{frame}{Exposure and distance decay}
\small
\[
w_{ij}(h)=
\frac{K(d_{ij}/h)}
{\sum_{\ell}K(d_{i\ell}/h)}
\]

\[
V_{ij}=w_j-r_i+A_i-c_{ij}+\varepsilon_{ij}
\]

\begin{itemize}
    \item A small bandwidth makes exposure local; a large bandwidth lets distant places matter.
    \item Travel costs may be Euclidean, network-based, flow-based, or generalized.
    \item Net value objects combine wages, rents, amenities, and access costs.
    \item The curation question is whether the measure matches the mechanism.
\end{itemize}
\end{frame}

\begin{frame}{Core operations}
\small
\begin{tabular}{p{0.23\linewidth} p{0.67\linewidth}}
\toprule
Operation & What to document \\
\midrule
Geocoding & provider, date, match score, failures, manual audit \\
Projection & input CRS, analysis CRS, distance/area rationale \\
Spatial join & within, intersects, centroid, nearest, ties, boundary cases \\
Buffer & radius, mechanism, alternative radii, barriers \\
Travel time & mode, departure time, routing source, missing routes \\
Crosswalk & source vintage, target geography, weight type, splits \\
Aggregation & denominator, weights, baseline year, support \\
Disclosure control & jitter, masking, suppression, precision loss \\
\bottomrule
\end{tabular}
\end{frame}

\begin{frame}{Boundary crosswalks}
\small
\[
X_h=\sum_g a_{gh}X_g
\]

\begin{itemize}
    \item \(a_{gh}\) can be area weight, population weight, employment weight, or flow weight.
    \item Area weights allocate land; population weights allocate residents.
    \item Employment weights allocate jobs; commuting-flow weights allocate market connections.
    \item A many-to-many crosswalk should be preserved before any collapsed file is used.
\end{itemize}
\end{frame}

\begin{frame}{Practical pitfalls}
\begin{itemize}
    \item MAUP: results change with aggregation unit.
    \item Geocoding error: failures and low-confidence matches are often systematic.
    \item Temporal mismatch: outcomes, boundaries, addresses, and routes come from different years.
    \item Endogenous geography choice: radii or units chosen after seeing results.
    \item Ecological fallacy: area averages interpreted as individual mechanisms.
    \item Edge effects: access or exposure just outside the study area is omitted.
    \item Support problems: some places have no comparable observations or valid routes.
\end{itemize}
\end{frame}

\begin{frame}{Recommended workflow}
\begin{enumerate}
    \item State the mechanism before choosing the map layer.
    \item Define unit of observation, exposure, and variation separately.
    \item Preserve raw spatial identifiers and source files.
    \item Validate geocoding, CRS, joins, crosswalks, and routes.
    \item Build the primary object and one plausible alternative.
    \item Audit edge cases, missing routes, split units, and out-of-support places.
    \item Document confidentiality masking and measurement consequences.
    \item Make the workflow rerunnable without manual GIS clicks.
\end{enumerate}
\end{frame}

\begin{frame}{Applied paper anchors}
\begin{itemize}
    \item Autor, Dorn, Hanson: commuting zones define local labor-market exposure.
    \item Chetty, Hendren, Katz: tract-level childhood exposure and timing matter.
    \item Monte, Redding, Rossi-Hansberg: commuting and migration connect local markets.
    \item Miller: job suburbanization makes access depend on worker and job locations.
    \item Donaldson, Hornbeck: market access depends on networks and trade costs.
\end{itemize}

\medskip
The lesson is not ``use GIS.'' The lesson is: define the economic object.
\end{frame}

\begin{frame}{Research Lab design}
\begin{columns}[T,onlytextwidth]
\begin{column}{0.32\linewidth}
\textbf{Reproduce}
\begin{itemize}
    \item synthetic tracts
    \item commuting zones
    \item industry shocks
    \item baseline weights
    \item exposure object
\end{itemize}
\end{column}
\begin{column}{0.32\linewidth}
\textbf{Diagnose}
\begin{itemize}
    \item county vs CZ
    \item crosswalk weights
    \item geocoding jitter
    \item edge cases
    \item estimand memo
\end{itemize}
\end{column}
\begin{column}{0.32\linewidth}
\textbf{Transfer}
\begin{itemize}
    \item job centers
    \item buffers
    \item distance decay
    \item travel time
    \item access memo
\end{itemize}
\end{column}
\end{columns}
\end{frame}

\begin{frame}{Where curation stops}
\begin{itemize}
    \item Curation defines places, assignments, exposure, distance, and access.
    \item Curation can diagnose measurement sensitivity and support.
    \item It does not prove exogeneity, solve sorting, or remove spillovers.
    \item Spatial causal inference begins when the researcher claims a counterfactual effect.
    \item Lecture 16 takes up spillovers, interference, border designs, and spatial dependence.
\end{itemize}
\end{frame}

\begin{frame}{Takeaways}
\begin{itemize}
    \item Geography is part of the estimand.
    \item Exposure weights are economic assumptions.
    \item Crosswalks, joins, buffers, and travel times need validation and documentation.
    \item Sensitivity to plausible geographies is a measurement diagnostic.
    \item Good spatial work is reproducible, mechanism-driven, and clear about confidentiality.
    \item Strong spatial curation makes the next lecture's causal designs possible.
\end{itemize}
\end{frame}

\end{document}
